\section{Direcciones futuras y desafíos de gobernanza}
\subsection{Entrenamiento de agentes y datos}
Graham señala que el entrenamiento actual de agentes todavía no aprovecha suficientemente los datos de «estilo agente» (agent chains, autorretroalimentación, corrección de fallas); en el futuro se necesitarán más agent-structured corpora para que el modelo aprenda a planificar y a corregir. Otra dirección es reforzar la evaluación human-in-the-loop: mediante una revisión conjunta entre humanos y agentes, el modelo aprende a detectar y corregir sus propios errores.

Agrega que el Agent training debería incorporar los registros del proceso como parte de los datos de entrenamiento, para que el modelo no solo vea el patch final, sino también el Plan, los pasos de ejecución y los resultados de Verify.

\subsection{Gobernanza y seguridad}
\begin{knowledgebox}{Núcleo de la investigación futura}
El énfasis está en: entrenar el comportamiento integrado del agente, incorporar trayectorias de corrección humana y construir evaluaciones más granulares (por ejemplo, si los registros del proceso están alineados y si el patch es mantenible).
\end{knowledgebox}

\begin{warningbox}{Gobernanza y seguridad}
La automatización excesiva de los agentes puede provocar decisiones de alto riesgo: por ejemplo, desplegar un patch directamente sin supervisión humana o invocar APIs sensibles entre repositorios; para llevarlo a producción es necesario introducir mecanismos de aprobación, auditoría y reversión.
\end{warningbox}

\begin{importantbox}{Desglose de direcciones}
Las tres rutas que recomienda Graham: 1) entrenar con agent-specific data, 2) construir o ampliar evaluaciones del tipo SWE-bench, 3) incorporar retroalimentación human-in-the-loop en tiempo real y mecanismos de seguridad.
\end{importantbox}

\subsection{Resumen de la sección}
\begin{itemize}
    \item Resolver el vacío en el entrenamiento y la evaluación de agentes es el punto de partida para que los agentes lleguen a una etapa de madurez.
    \item La seguridad, la aprobación y la intervención humana no deben considerarse añadidos, sino parte integral del ciclo de vida del agente.
\end{itemize}
